% !TEX program = xelatex
\documentclass[UTF8,fontset=fandol,a4paper,11pt]{ctexart}
\usepackage{../../_manual_common/hysim_manual}
\title{\textbf{HySim-XJTU-HRPES}\\[0.4em]{\Large validation 模块静态校验手册}\\
{\large 谓词、严重度、层级过滤与复杂度}}
\author{西安交通大学 HRPES 团队}
\date{\today}
\begin{document}
\maketitle
\begin{abstract}
本文档逐式对应 \srcpath{src/validation/validate\_system.cpp}、
\srcpath{include/hacdcpf/validation/validate\_system.hpp} 与 \srcpath{validation\_report.hpp}。
校验器产生结构化问题列表，不修复输入，也不证明数值求解收敛。
\end{abstract}
\tableofcontents

\section{输入、输出与数据结构}
输入为只读 \texttt{HybridPowerSystem} 和可选 \texttt{ValidationLevel}。输出
$R=[q_1,\ldots,q_m]$，每个问题
\begin{equation}
q=(severity,component\_type,component\_id,field,message).
\end{equation}
严重度枚举实际包含 \texttt{Info}、\texttt{Warning}、\texttt{Error}
（\srcpath{include/hacdcpf/validation/validation\_report.hpp:Severity}）。\texttt{count(s)} 在代码中统计“恰好等于” $s$
（:ValidationReport），并非注释所说的“at or above”；本文以函数体为准。

\section{报告有效性谓词}
\srcpath{ValidationReport::is\_valid}（\srcpath{include/hacdcpf/validation/validation\_report.hpp:ValidationReport}）的精确谓词为
\begin{equation}
valid(R)\Longleftrightarrow\neg\exists q\in R:q.severity=Error.
\end{equation}
因此只有 Warning 的报告仍然有效；$\texttt{has\_errors}=\neg valid(R)$，
$\texttt{has\_warnings}=\exists q:q.severity=Warning$，$\texttt{ok}=valid(R)$。Info 不影响这三个返回值。

\section{校验层级的精确过滤}
先令 $R_0=validate(sys)$。四级返回值由
\srcpath{src/validation/validate\_system.cpp:validate} 的 switch 直接定义：
\begin{align}
R_{SolverReady}&=R_0,\\
R_{Strict}&=\{promote(q):q\in R_0\},\\
promote(q).severity&=\begin{cases}Error,&q.severity=Warning,\\q.severity,&\text{否则，}\end{cases}\\
R_{Basic}&=\{q\in R_0:q.severity=Error\},\\
R_{Electrical}&=\{q\in R_0:q.severity=Error\lor
q.component\_type\in\{ACSystem,DCSystem\}\}.
\end{align}
因此 \texttt{Electrical} 并不是保留所有“拓扑相关” Warning，而是只按组件类型字符串保留
\texttt{ACSystem}/\texttt{DCSystem} Warning；ACBus 孤立负荷 Warning 会被过滤。这是当前代码行为，本文
不美化为更宽的语义。

\section{结构和数值谓词}
完整规则以 \srcpath{validate(sys)} 的顺序执行。代表性 helper 位于
\srcpath{src/validation/validate\_system.cpp:require\_order（lambda）}，与代码严格等价：
\begin{align}
require\_order(lo,hi)&\Rightarrow Error\quad\text{当且仅当 }lo>hi,\\
require\_nonnegative(x)&\Rightarrow Error\quad\text{当且仅当 }x<0,\\
require\_unit\_interval(x)&\Rightarrow Error\quad\text{当且仅当 }x<0\lor x>1,\\
require\_efficiency(\eta)&\Rightarrow Error\quad\text{当且仅当 }\eta\le0\lor\eta>1.
\end{align}
母线 ID 在 AC 与 DC 各自集合内检查正数和唯一性；跨域同号不报重复。支路端点必须不同且存在于本域。
AC/DC 母线均要求 $V_{min}>0$ 且 $V_{max}>V_{min}$。子系统与全局
\fld{base\_mva} 的差值仅在 $|S_{sub}-S_{sys}|>0.1$ 时产生 Warning。

\section{参考节点资格}
AC 参考母线支持谓词是以下任一在役设备成立：外部电网、发电机、构网且可控且
$P_{max}>P_{min}$ 的储能、控制角色提供 AC 角参考的 VSC。无支持时是 Error。DC 电压参考支持谓词是：
可控且 $P_{max}>P_{min}$ 的 DC 储能、提供 DC 电压参考或 PQ 模式的 VSC、输出端为该母线且处于 Voltage
或非零 droop 模式的可控 DC/DC、ConstantGamma/ConstantAlpha LCC。DC 无物理支持只产生 Warning，因为
PF 可把它作为理想双向电压边界。

\section{拓扑、算法流程与复杂度}
完整校验先扫描容器与引用，再调用 \texttt{build\_power\_system\_graph} 和
\texttt{analyze\_topology}。孤立负荷岛为每个 bus 添加 ACBus Warning；无 slack 的 AC 岛添加一个
ACSystem Warning。哈希集合建立和引用检查平均 $O(C)$，图分析为 $O(N+E)$；参考设备资格当前对每个参考
母线扫描多个设备容器，最坏 $O(N_{ref}C)$。

\section{异常、接口和验证建议}
校验器只追加问题，不抛出修复动作；NaN 的比较行为遵循 C++ 浮点比较，未被每个 helper 统一捕获，不能
宣称所有非有限量必然拒绝。上游为 model/io，求解入口应在 SolverReady 或业务指定层级后决定是否继续。
测试应逐条覆盖阈值两侧、四级过滤、AC/DC 同号 ID、无参考设备、理想 \texttt{DC\_V} Warning、孤岛类型和
\texttt{count} 的精确计数语义。

\section{工业指标、范围与改进}
指标包括错误规则分支覆盖率、误报/漏报率、诊断稳定 ID 完整率和百万组件扫描耗时。本模块只验证静态结构；
潮流残差、KKT、动态一致初值、I/O round-trip 和标准合规必须由对应模块验证。后续应把规则表生成化，并
修正文档注释与 \texttt{count} 实现语义的不一致。
% theory_validation_predicates.tex -- 静态校验理论层
\section{静态模型校验的谓词、层级与可靠性边界}\label{sec:validation-theory}

\begin{theorynote}
本章把静态校验器视为从系统到诊断集合的纯函数。它验证结构、引用、边界和参考节点资格，
不证明 PF/OPF 收敛，也不替输入执行修复。
\end{theorynote}

\subsection{诊断集合与有效性}
令 $Q$ 为问题集合，严重度为 Info、Warning、Error，则
\begin{equation}
valid(Q)\iff\nexists q\in Q:\ severity(q)=Error.
\end{equation}
四级校验是同一基础问题集合的过滤/提升：Basic 保留 Error，Electrical 额外保留
ACSystem/DCSystem 相关诊断，SolverReady 保持自然严重度，Strict 将 Warning 提升为 Error。

\subsection{结构谓词}
代表性约束包括
\begin{align}
lo&\le hi, & x&\ge0, & 0<\eta&\le1,\\
V_{min}&>0, & V_{max}&>V_{min}, & b_{from}&\ne b_{to}.
\end{align}
AC/DC 母线 ID 在各自域内唯一；同号 AC/DC ID 不构成重复。支路端点必须存在于同域集合。

\subsection{与实现的对应}
\begin{center}\footnotesize
\begin{longtable}{@{}L{7.0cm}L{8.2cm}@{}}
\toprule 谓词/层级 & 实现状态\\\midrule
\endfirsthead\toprule 谓词/层级 & 实现状态\\\midrule\endhead
基础字段与引用检查 & \implfull{src/validation/validate\_system.cpp:validate}\\
参考节点资格 & \implfull{src/validation/validate\_system.cpp:validate\_reference\_bus\_eligibility}\\
四级过滤/Warning 提升 & \implfull{src/validation/validate\_system.cpp:validate}\\
非有限值的统一拒绝 & \implpart{部分比较依赖 C++ NaN 语义，未由每个 helper 统一捕获}\\
PF/OPF 收敛证明 & \implnone\\
\bottomrule
\end{longtable}
\end{center}

% numerical_cross_validation.tex -- validation 数值证据
\section{规则边界、层级过滤与回归结果}\label{sec:validation-numerical}

macOS Release 的 \srcpath{test_validation} 通过 37 cases/178 assertions，覆盖空系统、AC/DC
同号 ID、重复 ID、端点引用、电压/阻抗/容量边界、参考节点、孤岛和四级过滤。

\begin{center}\small
\begin{tabular}{@{}L{6.5cm}L{5.0cm}@{}}
\toprule 校核对象 & 实测证据\\\midrule
基础规则与错误严重度 & 178 assertions\\
AC/DC 同号 ID 分域 & 通过\\
Warning 过滤与 Strict 提升 & 通过\\
参考节点资格与孤岛 & 通过\\
NaN 全字段拒绝 & 未形成统一独立门槛\\
\bottomrule
\end{tabular}
\end{center}

这些是静态谓词回归，不是误报/漏报率或百万组件性能认证。下游 PF、OPF、I/O round-trip、
动态一致初值必须由各自模块继续验证；校验器返回 OK 不等于数值求解一定收敛。

\section{跨模块工业级技术评价基线}

本章规定本手册所述模块进入工程算例、批量研究或生产服务前必须共同满足的评价基线。
具体模型公式、变量、约束和专用算法以正文相应章节为准；本章不扩大源码已经实现的范围。

\subsection{输入、输出、边界条件与数据结构审计}
输入审计应逐项检查有限性、单位、上下界、枚举值和引用完整性。系统对象必须先按所需层级完成
校验；AC 与 DC 母线采用域限定映射，组件稳定 \texttt{index}、向量位置和临时图索引不得混用。
输出审计必须记录单位、索引空间、模型范围、收敛或可行状态、后端、容差、迭代/节点计数和告警。
空系统、空时域、孤岛、重复 ID、零容量、退化约束与不可用可选后端均属于边界测试集。

\subsection{工程假设与数值稳定性分析}
模型使用的平衡、线性化、稳态、独立性、概率分布、时间离散或网络等值假设必须与应用场景匹配。
数值评价至少包含变量缩放、绝对/相对残差、可行性违反、条件数或枢轴质量、步长/阻尼、迭代停滞
和随机种子复现性。若采用正则化、平滑、回退、松弛或近似，应同时报告触发条件、参数值、对主指标
的敏感性以及是否改变模型口径；不得用“已返回结果”替代收敛或有效性证明。

\subsection{复杂度与资源分析}
令 $n_x$、$n_c$、$n_t$、$n_s$、$|V|$、$|E|$ 分别表示变量、约束、时段、场景、节点和边数。
报告应分别给出建模、求解、后处理的时间与内存。图遍历通常为 $O(|V|+|E|)$；逐时段/逐场景
流程至少线性增长为 $O(n_t n_s)$ 次子问题调用；稠密线性代数最坏可达 $O(n_x^3)$，稀疏方法则应
报告结构非零数、填充率和因子分解峰值内存；MILP/NLP 不给出虚假的多项式最坏界，而报告节点数、
间隙、时限和实测扩展曲线。复杂度结论必须基于固定硬件、固定后端和可复现算例族。

\subsection{异常工况处理}
非法输入应在进入求解器前返回结构化错误；不可行、无界、不收敛、时限、内存不足和后端缺失必须
相互区分。部分结果只有在对应 \fld{ValidityFlags}、\fld{model\_scope}、
\fld{model\_limitations} 或等价字段允许时才可使用。异常路径不得静默丢弃 AC/DC 资产、制造平衡
节点、把 NaN 改写为零或把近似解标为全局最优；系统原始对象应保持可恢复和可重试。

\subsection{与其他模块的接口关系}
接口链路遵循“工程语义模型 $\to$ 校验 $\to$ 投影/装配 $\to$ 求解或分析 $\to$ 结果回投影”。
跨模块传递时保留稳定 ID、单位、符号、时间基准和场景身份；消费潮流、OPF、动态或可靠性结果前，
先检查其收敛、覆盖范围和有效性标志。HTTP/JSON/Python 层只做语义保持适配，不重新解释求解结果。

\subsection{测试与验证建议}
最低验证矩阵包括：手算小例、标准算例、边界/异常输入、随机性质测试、算法或后端交叉校核、
序列化往返、稳定 ID 回投影、AC/DC 同号母线、重复运行确定性和规模扩展测试。数值算法还应加入
残差独立重算、雅可比有限差分或自动微分校核；近似模型应与高保真模型比较并给出误差分布，而非
只报告单个平均值。回归报告必须列明构建配置、算例版本、容差、通过/失败/跳过数及未验证范围。

\subsection{工业级评估指标}
\begin{itemize}
  \item \textbf{正确性}：守恒残差、约束最大违反、解析/基准误差、结果回投影一致性；
  \item \textbf{鲁棒性}：收敛率、不可行识别率、异常分类准确率、扰动敏感性与随机复现率；
  \item \textbf{性能}：端到端时延、吞吐量、峰值内存、并行效率与规模增长斜率；
  \item \textbf{可追溯性}：输入模式版本、稳定 ID 覆盖率、模型范围/限制字段完整率、告警可定位率；
  \item \textbf{工程有效性}：与标准、外部引擎或现场数据的偏差，以及误判/漏判对业务指标的影响。
\end{itemize}
指标应预先给出验收阈值和失败处置；未达到阈值时只能标记为实验性或受限适用。

\subsection{适用范围、局限性与后续改进}
本手册只评价当前源码覆盖的模型、后端和数据结构，不对未建模物理过程、未安装求解器或未执行的
外部验证作保证。后续改进应按“缺口 $\to$ 可量化指标 $\to$ 源码/测试变更 $\to$ 独立验证”闭环，
优先消除会造成错误结论的语义缺口，其次提升数值鲁棒性与性能，最后扩展模型范围。

\end{document}
